\section{占用率、资源约束与线程块大小}

\subsection{占用率的定义}

\term{占用率}{occupancy}描述一个流式多处理器上实际驻留的线程束数量，相对于该流式多处理器允许的最大驻留线程束数量的比例。记实际驻留线程束数为 \(W_{\mathrm{res}}\)，最大驻留线程束数为 \(W_{\max}\)，则
\[
\boxed{\occupancy=\frac{W_{\mathrm{res}}}{W_{\max}}}.
\]

由于同一架构上的线程束大小固定，若只考虑线程数约束，也可以等价地用驻留线程数与最大驻留线程数之比理解。但正式定义以驻留线程束为中心更直接，因为调度器的候选单位就是线程束。

\begin{warningbox}
占用率中的“驻留/活动”不等于“当前周期正在执行”。一个线程束可以已经驻留，却正在等待内存；也可以已经就绪，却暂时没有得到发射机会。占用率衡量的是可驻留并保持状态的工作池大小。
\end{warningbox}

\subsection{流式多处理器资源示例}

不同 GPU 架构具有不同的资源参数。最大驻留线程数通常显著高于每个流式多处理器的核心数，同时还受到线程块数、寄存器和共享内存等约束。

\begin{figure}[H]
  \centering
  \includegraphics[width=\textwidth]{sm_resources_p27.png}
  \caption{多代 GPU 的流式多处理器资源对比；红框突出每个流式多处理器的核心数与最大驻留线程数}
\end{figure}

\begin{table}[H]
\centering
\small
\setlength{\tabcolsep}{4pt}
\begin{tabularx}{\textwidth}{>{\raggedright\arraybackslash}Xrrrrrr}
\toprule
GPU & K40 & M40 & P100 & V100 & A100 & H100 \\
\midrule
微架构 & Kepler & Maxwell & Pascal & Volta & Ampere & Hopper \\
计算能力 & 3.5 & 5.2 & 6.0 & 7.0 & 8.0 & 9.0 \\
流式多处理器数量 & 15 & 24 & 56 & 80 & 108 & 132 \\
每个流式多处理器核心数 & 192 & 128 & 64 & 64 & 64 & 128 \\
最大驻留线程数/流式多处理器 & 2048 & 2048 & 2048 & 2048 & 2048 & 2048 \\
最大驻留线程块数/流式多处理器 & 16 & 32 & 32 & 32 & 32 & 32 \\
最大线程数/线程块 & 1024 & 1024 & 1024 & 1024 & 1024 & 1024 \\
寄存器/流式多处理器 & 64K & 64K & 64K & 64K & 64K & 64K \\
共享内存/流式多处理器 & 48KB & 96KB & 64KB & 96KB & 164KB & 228KB \\
\bottomrule
\end{tabularx}
\caption{多代 GPU 的流式多处理器资源示例}
\end{table}

\subsection{把多重驻留限制写成统一形式}

设一个线程块有 \(T_b\) 个线程，线程束大小为 \(W\)，则每个线程块需要
\[
W_b=\left\lceil\frac{T_b}{W}\right\rceil
\]
个线程束。设设备限制分别为：最大驻留线程块数 \(B_{\max}\)、最大驻留线程数 \(T_{\max}\)、最大驻留线程束数 \(W_{\max}\)。若每个线程块还需要 \(R_b\) 个寄存器资源单位与 \(S_b\) 字节共享内存，而流式多处理器总量分别为 \(R_{\mathrm{SM}}\) 和 \(S_{\mathrm{SM}}\)，则多重限制可概念化为
\[
B_{\mathrm{res}}
\le
\min\left(
B_{\max},
\left\lfloor\frac{T_{\max}}{T_b}\right\rfloor,
\left\lfloor\frac{W_{\max}}{W_b}\right\rfloor,
\left\lfloor\frac{R_{\mathrm{SM}}}{R_b}\right\rfloor,
\left\lfloor\frac{S_{\mathrm{SM}}}{S_b}\right\rfloor
\right).
\]

该表达式用于统一理解各类驻留约束；实际硬件还可能按特定粒度分配寄存器或共享内存，因此精确驻留数应由设备属性与占用率工具计算。确定 \(B_{\mathrm{res}}\) 后，驻留线程束数为
\[
W_{\mathrm{res}}=B_{\mathrm{res}}W_b,
\]
再代入占用率定义即可。

\subsection{H100 风格限制下的三个线程块大小}

以下例题采用一组简化限制：每线程块最多 1024 个线程，每流式多处理器最多 2048 个驻留线程、最多 32 个驻留线程块，并假设寄存器和共享内存不形成额外限制。

\subsubsection{每块 256 个线程}

线程数约束允许
\[
\left\lfloor\frac{2048}{256}\right\rfloor=8
\]
个线程块，且 \(8<32\)。所以可驻留 \(8\times256=2048\) 个线程，即 64 个线程束；在这一简化模型中占用率为
\[
\occupancy=\frac{64}{64}=100\%.
\]

\subsubsection{每块 32 个线程}

若只看线程数，需要
\[
2048/32=64
\]
个线程块才能填满线程容量，但设备最多只允许 32 个驻留线程块。因此实际只能驻留
\[
32\times32=1024
\]
个线程，也就是 32 个线程束，对应简化占用率
\[
\occupancy=\frac{32}{64}=50\%.
\]
这里的瓶颈是\textbf{最大驻留线程块数}。

\subsubsection{每块 768 个线程}

线程数约束允许
\[
\left\lfloor\frac{2048}{768}\right\rfloor=2
\]
个线程块，驻留线程数为
\[
2\times768=1536.
\]
剩余 \(2048-1536=512\) 个线程容量无法再容纳完整的第三个线程块，因此简化占用率为
\[
\occupancy=\frac{1536}{2048}=75\%.
\]

\begin{keybox}
线程块驻留必须以完整线程块为单位，所以资源利用会出现“装箱”效应。即使总资源还有余量，只要余量不足以容纳另一个完整线程块，这部分资源就不能被该核函数的额外线程块使用。
\end{keybox}

\subsection{二维线程块大小：彩色转灰度核函数}

另一个例题假设某 GPU 的每个流式多处理器最多驻留 1536 个线程和 8 个线程块，比较 \(8\times8\)、\(16\times16\) 与 \(32\times32\) 三种二维线程块。

\begin{center}
\begin{tabular}{lrrrr}
\toprule
线程块形状 & 每块线程数 & 可驻留块数 & 驻留线程数 & 简化占用率 \\
\midrule
\(8\times8\) & 64 & 8 & 512 & \(33.3\%\) \\
\(16\times16\) & 256 & 6 & 1536 & \(100\%\) \\
\(32\times32\) & 1024 & 1 & 1024 & \(66.7\%\) \\
\bottomrule
\end{tabular}
\end{center}

推导过程如下：

\begin{itemize}
  \item \(8\times8\)：线程数允许 \(1536/64=24\) 个线程块，但块数上限为 8，所以只能驻留 512 个线程；
  \item \(16\times16\)：\(1536/256=6\) 个线程块，既不超过 8 块限制，又正好使用全部 1536 个线程容量；
  \item \(32\times32\)：每块 1024 线程，第二块会使线程数达到 2048，因此只能驻留一块。
\end{itemize}

在该题给定的简化限制下，\(16\times16\) 获得最高理论占用率。

\subsection{占用率高并不保证性能最高}

较高占用率通常有利于增加延迟隐藏机会，但占用率只在能够帮助隐藏延迟的范围内有用，并不单独保证性能。

原因与上一节一致：若内存带宽已经饱和、所有线程束都受共同依赖阻塞，或其他执行资源成为瓶颈，即使继续增加驻留线程束，也未必增加有效指令吞吐。因此，占用率最适合作为\textbf{资源驻留与延迟隐藏能力的诊断量}，而不是完整的性能模型。
